% ============================================================
% CAPITOLO 14 — Testing e Sicurezza
% ============================================================
\chapter{Testing e Sicurezza}\index{testing}\index{sicurezza}

\section*{Cosa Costruirai}
Un sistema di testing e verifica di sicurezza per l'intero stack:
\begin{itemize}
  \item Test unitari generati dall'IA per backend e frontend
  \item Test di integrazione per gli endpoint API
  \item Analisi di sicurezza OWASP Top~10 automatizzata
  \item Confidence Tagging per decidere cosa revisionare manualmente
  \item Pipeline di qualità riproducibile
\end{itemize}

\textbf{Tempo stimato}: 60--90 minuti\\
\textbf{Prerequisito}: Applicazione full-stack completa (Capitoli~6--13)

\begin{dangerbox}[title={La Sicurezza Non è Opzionale}]
Il codice generato dall'IA può contenere vulnerabilità di sicurezza --- non per malevolenza, ma perché l'IA ottimizza per la funzionalità, non per la sicurezza, a meno che non glielo chiedi esplicitamente. Questo capitolo ti fornisce gli strumenti per individuare i problemi più comuni (OWASP Top~10), ma \textbf{non sostituisce una revisione professionale di sicurezza}.

\textbf{Prima di un deployment commerciale} --- un'applicazione che gestisce dati di clienti reali, transazioni finanziarie, dati sanitari o qualsiasi informazione sensibile --- fai revisionare il codice da un professionista di sicurezza informatica. Il costo di un audit è una frazione del danno reputazionale e legale di un data breach.

Le tecniche che impari qui sono il tuo \textbf{primo livello di difesa}, non l'ultimo.
\end{dangerbox}

% ---------------------------------------------------------
\section{Test del Backend}

\subsection{Il principio: Test come Documentazione}

I test generati dall'IA hanno un doppio valore: verificano il funzionamento \textbf{e} documentano il comportamento atteso. Se domani cambi qualcosa e un test fallisce, il messaggio di errore ti dice esattamente cosa si è rotto.

\begin{praticabox}[title={Pratica --- Genera i test del backend}]
\begin{lstlisting}[language={}]
Genera test completi per il backend Notes API con vitest.

Setup:
1. Installa vitest e supertest come devDependencies
2. Crea vitest.config.js con database PostgreSQL di test
3. Esegui migrazioni prima dei test, pulisci i dati tra i test

Test richiesti:

test/unit/
  - note-service.test.js: CRUD del servizio note
    (crea, leggi, aggiorna, elimina, nota non trovata, accesso negato)
  - auth-service.test.js: token (generazione JWT, verifica,
    refresh, token scaduto)
  - validation.test.js: middleware Zod (titolo vuoto,
    contenuto troppo lungo, campi extra ignorati)

test/integration/
  - notes-api.test.js: endpoint con supertest
    (GET /api/notes senza auth -> 401,
     POST con dati validi -> 201,
     GET /:id di altro utente -> 403,
     DELETE /:id -> 200 + nota rimossa)
  - auth-api.test.js: flusso auth
    (GET /api/auth/me senza token -> 401,
     con token valido -> 200 + profilo,
     POST /api/auth/refresh con refresh valido -> nuovo access)

Ogni test deve avere un nome descrittivo che spiega lo scenario.
\end{lstlisting}
\end{praticabox}

\begin{lstlisting}[language=bash,title={Esegui i test}]
cd backend
npx vitest run
\end{lstlisting}

Output atteso:
\begin{lstlisting}[language={}]
 v note-service.test.js (6 tests)
 v auth-service.test.js (4 tests)
 v validation.test.js (3 tests)
 v notes-api.test.js (4 tests)
 v auth-api.test.js (3 tests)

 Test Files  5 passed
 Tests       20 passed
\end{lstlisting}

\begin{tipbox}
Se un test fallisce, copia l'errore e incollalo a Copilot: ``Questo test fallisce con questo errore. Qual è il problema?'' L'IA di solito identifica il bug al primo tentativo.
\end{tipbox}

% ---------------------------------------------------------
\section{Test del Frontend}

\begin{praticabox}[title={Pratica --- Genera i test React}]
\begin{lstlisting}[language={}]
Genera test per il frontend React con vitest e @testing-library/react.

1. Installa: vitest, @testing-library/react,
   @testing-library/jest-dom, jsdom, msw (Mock Service Worker)
2. Configura MSW per simulare il backend

Test richiesti:
test/components/
  - NoteCard.test.jsx: renderizza titolo e contenuto
  - NoteForm.test.jsx: validazione campi, submit callback
  - ProtectedRoute.test.jsx: redirige se non autenticato

test/hooks/
  - useNotes.test.jsx: carica note, gestisce errori, aggiorna

test/pages/
  - LoginPage.test.jsx: mostra bottoni OAuth
  - DashboardPage.test.jsx: loading, lista note, empty state
\end{lstlisting}
\end{praticabox}

\begin{lstlisting}[language=bash]
cd frontend
npx vitest run
\end{lstlisting}

% ---------------------------------------------------------
\section{Test dell'App Flutter}

\begin{praticabox}[title={Pratica --- Genera i test Flutter}]
\begin{lstlisting}[language={}]
Genera test per l'app Flutter:

test/
  unit/
    - token_service_test.dart: salvataggio, lettura, cancellazione
    - note_model_test.dart: fromJson, toJson, campi obbligatori

  widget/
    - login_screen_test.dart: bottoni OAuth, loading state
    - note_card_test.dart: renderizza titolo, tap callback
    - dashboard_screen_test.dart: loading -> lista note, empty state

Usa mockito per simulare i servizi.
Usa ProviderScope.overrides per iniettare provider mock.
\end{lstlisting}
\end{praticabox}

\begin{lstlisting}[language=bash]
cd notes_mobile
flutter test
\end{lstlisting}

% ---------------------------------------------------------
\section{Confidence Tagging}\index{Confidence Tagging}

Non tutto il codice generato dall'IA richiede la stessa attenzione in revisione. Il \textbf{Confidence Tagging} classifica il codice per livello di rischio.

\subsection{La matrice di rischio}

\begin{center}
\begin{tabularx}{\textwidth}{l l X}
\toprule
\textbf{Livello} & \textbf{Cosa include} & \textbf{Azione} \\
\midrule
\textcolor{green!60!black}{\textbf{LOW}} & UI, styling, componenti stateless, modelli dati & Revisione rapida: il codice funziona? \\
\textcolor{orange}{\textbf{MEDIUM}} & Logica business, validazione, routing, state management & Revisione attenta: la logica è corretta? \\
\textcolor{red}{\textbf{HIGH}} & Auth, autorizzazione, crypto, query DB, input utente & Revisione rigorosa: è sicuro? \\
\bottomrule
\end{tabularx}
\end{center}

\begin{praticabox}[title={Pratica --- Analisi di rischio del progetto}]
\begin{lstlisting}[language={}]
Analizza l'intero progetto e classifica ogni file:

- LOW RISK: non gestiscono dati sensibili ne logica critica
- MEDIUM RISK: logica business o validazione
- HIGH RISK: autenticazione, autorizzazione, dati utente, query DB

Per ogni file HIGH RISK, elenca le vulnerabilita da verificare:

HIGH: backend/src/middleware/auth.js
   - JWT: verificare che il secret non sia hardcoded
   - Token: gestisca token malformati
   - User lookup: non vulnerabile a injection

MEDIUM: backend/src/services/note-service.js
   - Verificare che filtri sempre per userId

LOW: frontend/src/components/ui/Button.jsx
   - Nessuna verifica necessaria
\end{lstlisting}
\end{praticabox}

% ---------------------------------------------------------
\section{Code Review del Codice Generato dall'IA}\index{code review}

Il Confidence Tagging ti dice \textbf{cosa} revisionare. Questa sezione ti insegna \textbf{come} farlo --- anche se non sei uno sviluppatore esperto.

\subsection{Perché l'IA sbaglia in modi diversi dall'umano}

Un programmatore umano fa errori di distrazione: typo, parentesi dimenticate, variabili con nomi sbagliati. L'IA non fa quasi mai questi errori. I bug dell'IA sono più insidiosi:

\begin{center}
\begin{small}
\begin{tabularx}{\textwidth}{l l X}
\toprule
\textbf{Tipo di bug} & \textbf{Esempio} & \textbf{Perché succede} \\
\midrule
Happy path only & Funziona con input validi, crasha con input vuoti & L'IA ottimizza per ``funziona'', non per ``resiste'' \\
API inesistenti & Usa un metodo che non esiste nella versione installata & Ha mescolato documentazione di versioni diverse \\
Logica plausibile ma errata & Un filtro che sembra corretto ma esclude casi limite & Il pattern sembrava giusto nel training data \\
Sicurezza dimenticata & Endpoint senza controllo di autorizzazione & Non glielo hai chiesto esplicitamente \\
Dipendenze fantasma & Importa un pacchetto non nel package.json & L'ha visto in un progetto simile \\
\bottomrule
\end{tabularx}
\end{small}
\end{center}

\subsection{Metodologia di review per livello}

\textbf{\textcolor{green!60!black}{LOW RISK} --- Verifica funzionale (2 minuti/file)}

Componenti UI, styling, modelli dati. La domanda è una sola: \textit{funziona?}

\begin{enumerate}
  \item L'IA ha creato il file? Aprilo e verifica che non sia vuoto
  \item Avvia l'app e guarda se il componente si visualizza
  \item Se sì, passa oltre
\end{enumerate}

\textbf{\textcolor{orange}{MEDIUM RISK} --- Verifica logica (5--10 minuti/file)}

Logica business, validazione, routing. La domanda: \textit{fa quello che dovrebbe?}

\begin{enumerate}
  \item Leggi la funzione principale e chiediti: ``Cosa succede se l'input è vuoto? Nullo? Enorme?''
  \item Cerca le parole chiave pericolose: \texttt{any}, \texttt{as any} (TypeScript), \texttt{!} (force unwrap), \texttt{catch} vuoti
  \item Verifica che le validazioni Zod/Joi corrispondano ai requisiti del \texttt{\_CONTEXT.md}
  \item Se la logica ti sembra corretta ma non sei sicuro, chiedi all'IA: \textit{``Spiega passo per passo cosa fa questa funzione e quali edge case potrebbe non gestire''}
\end{enumerate}

\textbf{\textcolor{red}{HIGH RISK} --- Verifica rigorosa (15--30 minuti/file)}

Autenticazione, autorizzazione, crypto, query DB. La domanda: \textit{è sicuro?}

\begin{enumerate}
  \item \textbf{Autorizzazione}: ogni endpoint che accede a dati utente filtra per \texttt{userId}? Può un utente accedere ai dati di un altro?
  \item \textbf{JWT}: il secret è in una variabile d'ambiente? Il token ha una scadenza?
  \item \textbf{Input utente}: tutto l'input passa attraverso validazione prima di raggiungere il database?
  \item \textbf{Errori}: i messaggi di errore in produzione \textsc{non} espongono stack trace o dettagli interni?
  \item Usa la Security Checklist come guida sistematica
\end{enumerate}

\begin{praticabox}[title={Pratica --- Review assistita dall'IA}]
\begin{lstlisting}[language={}]
Analizza i file HIGH RISK del progetto notes-fullstack.

Per ogni file:
1. Elenca tutti i possibili edge case non gestiti
2. Identifica ogni punto dove l'input utente raggiunge il database
   senza passare per validazione
3. Verifica che non ci siano segreti hardcoded
4. Controlla che ogni endpoint protetto verifichi l'autorizzazione

Formato output per ogni file:
VERIFICATO: [cosa va bene]
ATTENZIONE: [cosa potrebbe essere un problema]
BUG: [cosa e sicuramente sbagliato, con fix]
\end{lstlisting}
\end{praticabox}

\begin{tipbox}
\textbf{La regola d'oro della review}: Non devi capire \textit{come} funziona ogni riga di codice. Devi capire \textit{cosa} fa la funzione e chiederti: ``Cosa potrebbe andare storto?''. L'IA è il tuo microscopio --- tu devi sapere dove puntarlo.
\end{tipbox}

% ---------------------------------------------------------
\section{Analisi OWASP Top 10}\index{OWASP!Top 10}

Le 10 vulnerabilità più comuni nelle applicazioni web~\cite{owasp:top10}. Verifichiamole nel nostro progetto.

\begin{praticabox}[title={Pratica --- Audit OWASP}]
\begin{lstlisting}[language={}]
Esegui un'analisi di sicurezza basata su OWASP Top 10 (2021).

A01: Broken Access Control
- Un utente puo accedere alle note di un altro?
- Gli endpoint admin sono protetti?
- Le route protette verificano SEMPRE il JWT?

A02: Cryptographic Failures
- Il JWT secret e sufficientemente lungo?
- Le password sono hashate con bcrypt?
- Il database usa SSL in produzione?

A03: Injection
- Le query Prisma sono parametrizzate?
- L'input e validato con Zod prima del database?

A04: Insecure Design
- Il refresh token ha una scadenza?
- C'e rate limiting sul login?

A05: Security Misconfiguration
- Errori in produzione espongono dettagli interni?
- Helmet.js configurato?
- CORS configurato per la produzione?

A06: Vulnerable Components
- npm audit non riporta vulnerabilita critiche?

A07: Authentication Failures
- I token JWT scadono?
- Il logout invalida il refresh token nel database?

A08: Software and Data Integrity
- package-lock.json committato?

A09: Security Logging
- I tentativi di login falliti vengono loggati?
- I log NON contengono dati sensibili?

A10: SSRF
- Il backend fa richieste a URL forniti dall'utente?

Per ogni vulnerabilita trovata, fornisci il fix con il codice.
\end{lstlisting}
\end{praticabox}

\subsection{Fix comuni}

\textbf{Rate Limiting:}\index{rate limiting}
\begin{lstlisting}[language={}]
Aggiungi rate limiting al backend:
1. Installa express-rate-limit
2. Limita /api/auth/* a 10 richieste/minuto per IP
3. Limita /api/notes a 100 richieste/minuto per utente
\end{lstlisting}

\textbf{Helmet.js:}\index{Helmet.js}
\begin{lstlisting}[language={}]
Aggiungi Helmet.js per gli header di sicurezza:
1. Installa helmet
2. Aggiungi helmet() come primo middleware
3. Configura CSP per la produzione
\end{lstlisting}

\textbf{Audit dipendenze:}
\begin{lstlisting}[language=bash]
cd backend && npm audit
cd frontend && npm audit
npm audit fix  # se vengono riportate vulnerabilita
\end{lstlisting}

% ---------------------------------------------------------
\section{Security Checklist Finale}

\begin{praticabox}[title={Pratica --- Verifica completa}]
\begin{center}
\begin{small}
\begin{tabularx}{\textwidth}{c l X}
\toprule
\textbf{\#} & \textbf{Controllo} & \textbf{Stato} \\
\midrule
1 & JWT secret $\geq$ 256 bit, in variabile d'ambiente & $\square$ \\
2 & Access token scade in 1 ora & $\square$ \\
3 & Refresh token scade in 7 giorni & $\square$ \\
4 & Logout cancella il refresh token dal database & $\square$ \\
5 & Note filtrate per userId in \textsc{tutte} le query & $\square$ \\
6 & Input validato con Zod prima del database & $\square$ \\
7 & Errori in produzione \textsc{non} espongono stack trace & $\square$ \\
8 & Helmet.js attivo con header di sicurezza & $\square$ \\
9 & Rate limiting su endpoint auth & $\square$ \\
10 & CORS configurato solo per domini consentiti & $\square$ \\
11 & npm audit senza vulnerabilità critiche & $\square$ \\
12 & Token mobile in flutter\_secure\_storage & $\square$ \\
13 & Frontend non salva token in localStorage & $\square$ \\
14 & Database usa SSL in produzione & $\square$ \\
15 & \texttt{.env} e \texttt{keystore.jks} nel \texttt{.gitignore} & $\square$ \\
\bottomrule
\end{tabularx}
\end{small}
\end{center}
\end{praticabox}

\begin{checkpointbox}
Se tutti i 15 controlli sono spuntati e i test passano, l'applicazione è pronta per il deploy in produzione.
\end{checkpointbox}

% ---------------------------------------------------------
\section{Sicurezza Agentica: Oltre l'OWASP Web}

L'analisi OWASP Top~10 protegge l'\textbf{applicazione} che costruisci. Ma chi protegge l'\textbf{infrastruttura di sviluppo}? Quando concedi a un agente IA l'accesso a database, file system e API tramite MCP, introduci vettori di minaccia nuovi che l'OWASP tradizionale non copre.

Nel 2025 l'OWASP ha pubblicato la \textbf{OWASP Top~10 for MCP}\index{OWASP!MCP Top 10}~\cite{owasp:mcp-top10,owasp:ai-agent}, una classificazione specifica per le vulnerabilità del protocollo. Tre di queste meritano attenzione immediata:

\subsection{Tool Poisoning (MCP03)}\index{Tool Poisoning}\index{MCP!Tool Poisoning}

Quando configuri un server MCP, l'agente si fida delle \textbf{descrizioni} dei tool esposti dal server per decidere quale azione eseguire. Un server compromesso può inserire descrizioni malevole che inducono l'agente a eseguire operazioni distruttive.

\begin{dangerbox}
Un server MCP PostgreSQL manomesso potrebbe descrivere il tool \texttt{cleanup\_old\_data} come ``rimuove i record di test'', quando in realtà esegue \texttt{DROP TABLE users CASCADE}. L'agente, fidandosi della descrizione, lo invocherebbe senza esitazione.
\end{dangerbox}

\textbf{Mitigazione:} installa \textit{solo} server MCP da repository verificati (\texttt{@modelcontextprotocol/} ufficiali o con migliaia di stelle GitHub). Verifica il codice sorgente del server prima dell'adozione.

\subsection{Shadow MCP Servers (MCP09)}\index{Shadow MCP Server}

Come accadeva con lo ``Shadow IT'', gli sviluppatori possono installare server MCP non verificati per accelerare l'integrazione con nuove API. Questi server bypassano i controlli di sicurezza del team e creano punti di ingresso occulti.

\textbf{Mitigazione:} mantieni un registro condiviso dei server MCP approvati (nel \texttt{\_CONTEXT.md} di progetto o nelle policy del team).

\subsection{Context Injection e Over-Sharing (MCP10)}\index{Context Injection}\index{MCP!Context Injection}

Quando la finestra di contesto è condivisa tra sessioni diverse, informazioni sensibili --- token API, segreti, dati utente --- possono fuoriuscire da un task all'altro.

\textbf{Mitigazione:} usa sempre file \texttt{.env} per i segreti, \textbf{mai} inline nelle configurazioni MCP. Verifica che il \texttt{.gitignore} includa tutti i file sensibili. In ambienti di team, considera la mutual~TLS per autenticare le connessioni MCP.

\begin{praticabox}[title={Pratica --- Audit di sicurezza agentica}]
\begin{lstlisting}[language={}]
Analizza la sicurezza della nostra infrastruttura agentica:

1. TOOL VERIFICATION:
   - Elenca tutti i server MCP configurati (.vscode/mcp.json)
   - Per ognuno: verifica la fonte, controlla le stelle GitHub
     e la data dell'ultimo commit
   - Segnala server non ufficiali o non manutenuti

2. SECRET MANAGEMENT:
   - Verifica che NESSUN segreto sia hardcoded in mcp.json
   - Tutti i segreti devono essere in variabili d'ambiente
   - Il file .env è nel .gitignore?

3. CONTEXT HYGIENE:
   - Il _CONTEXT.md contiene credenziali?
   - I log di sessione contengono dati sensibili?
   - Le configurazioni MCP espongono URL o token interni?
\end{lstlisting}
\end{praticabox}

% ---------------------------------------------------------
\section{Testing di Sistemi Non Deterministici}

I test unitari e di integrazione che hai scritto in questo capitolo validano \textbf{codice deterministico}: dato un input, l'output è sempre lo stesso. Ma nel 2026 molte applicazioni integrano agenti IA che producono risposte \textbf{probabilistiche} --- chatbot, motori di raccomandazione, pipeline di ragionamento autonomo.

In questi scenari, le asserzioni binarie pass/fail crollano: l'output ``corretto'' ha molte forme possibili.

\subsection{Strategie per il testing probabilistico}

\begin{center}
\begin{tabularx}{\textwidth}{l X}
\toprule
\textbf{Tecnica} & \textbf{Quando usarla} \\
\midrule
Valutazione semantica & L'output deve \textit{significare} la stessa cosa, non essere identico lettera per lettera \\
Golden dataset & Confronta le risposte con un set di riferimento curato da esperti umani \\
Evaluator Agent & Un secondo agente valuta la qualità dell'output del primo \\
Osservabilità (tracing) & Traccia l'\textit{albero decisionale} dell'agente per capire \textit{perché} ha dato quella risposta \\
Fuzzing probabilistico & Inietta variazioni nell'input per verificare la robustezza dell'output \\
\bottomrule
\end{tabularx}
\end{center}

\begin{tipbox}
Piattaforme come \textbf{Maxim~AI}~\cite{maxim:eval-2026} permettono di valutare sistemi multi-agente misurando l'accuratezza del ragionamento su interazioni multi-turno. Strumenti open-source come \textbf{Arize Phoenix}~\cite{arize:phoenix} (basato su OpenTelemetry) tracciano l'intero albero di esecuzione delle decisioni di un agente.
\end{tipbox}

\begin{notebox}
Per le applicazioni di questo libro (CRUD, REST~API, autenticazione) i test tradizionali sono perfettamente adeguati. Il testing probabilistico diventa necessario solo quando integri modelli IA \textit{all'interno} della tua applicazione --- ad esempio un chatbot o un motore di ricerca semantica.
\end{notebox}

% ---------------------------------------------------------
\section*{Riepilogo}

\begin{center}
\begin{tabularx}{\textwidth}{l X}
\toprule
\textbf{Aspetto} & \textbf{Dettaglio} \\
\midrule
Test Backend & vitest + supertest, unit + integration \\
Test Frontend & vitest + testing-library + MSW \\
Test Mobile & flutter test + mockito \\
Rischio & Confidence Tagging a 3 livelli \\
Code Review & Metodologia di review per livello di rischio \\
Sicurezza Web & OWASP Top~10 audit con fix \\
Sicurezza Agentica & OWASP MCP Top~10: Tool Poisoning, Shadow MCP, Context Injection \\
Testing Probabilistico & Strategie per output non deterministici \\
Checklist & 15 controlli di sicurezza verificabili \\
\bottomrule
\end{tabularx}
\end{center}
